% Results fragment for Section 8 (Policy index and Strategic Index), subsection "Results".
% Generated numbers: paper/results/strategic.json; figures and table: scripts/paper/fig_strategic.R
% (final rerun of 2026-10-12; the script recomputes the Policy index, the Strategic Index, the
% endowment flag and the allied selection from the stored inputs with the pipeline's functions
% and stops unless all four equal the stored outputs and data/processed/charts/allied_strategic.csv).
% Figures: strategic_policy, strategic_heatmap, strategic_decomp, strategic_screen, strategic_sign.
% Table: tables/strategic_allied.tex (tab:strategic_allied).
% All correlations are Spearman rank correlations. Every Strategic Index cell requires rankable
% EO and ES cells (rankable == TRUE), so every ranking below is over rankable cells only.

\paragraph{Coverage of the Policy index.} The Policy index combines two components: the record of
national policies and measures in the IEA Policies and Measures database \citep{iea2026pams}
(weight 3), which is required, and the Climate Action Tracker rating \citep{cat2026ratings}
(weight 2), which is added where the Tracker rates the country; the industrial-policy measures of
the New Industrial Policy Observatory and the dual-use score are not part of the composite. The
IEA database covers 215 economies, and all 215 receive a Policy index
(Figure~\ref{fig:strategic_policy}, panel a). The Tracker rates 65 of them, 39 with a national
rating and 26 EU members through the rating of the Union, so 65 economies are scored on both
components and 150 on the IEA database alone; the basis is recorded for every country. All 57
coalition members have a Policy index: 46 on both components and 11 on the IEA database alone, two
of them core members (Israel and Ukraine) and nine partners (the Democratic Republic of the Congo,
the Dominican Republic, Gabon, Ghana, Guatemala, Malaysia, Namibia, Oman and Qatar). Under the
earlier rule that required both components, only the 65 economies with a Tracker rating would be
eligible, and these 11 members would have neither a Strategic Index nor an allied pick. The index
covers 6,762 cells; 203 economies have it in all 33 technology--stage cells and 12 territories in
some cells only.

The two components measure different things, and the composite inherits the difference
(Figure~\ref{fig:strategic_policy}, panel b). Over the 65 economies scored on both, the Spearman
rank correlation of the two components is $-0.05$ across cells and $-0.23$ across country means,
and within a country it has a median of $-0.02$. The Tracker's rating is one economy-wide grade
that varies across a country's cells only through the fixed technology greenhouse-gas index $g_t$
of the climate term: a rating of ``insufficient'' or better tilts the country's cells toward
low-emission technologies, ``highly insufficient'' leaves them untilted, and ``critically
insufficient'' tilts them toward high-emission technologies. Within a country the composite
therefore follows the Tracker component closely (median Spearman 0.79) and the IEA component only
partly (median 0.48, interquartile range 0.41--0.57). Ranking these countries' cells on the IEA
component alone instead of the composite gives a median within-country Spearman of 0.48 over all
their policy cells and 0.44 over their Strategic Index cells. The basis is thus not innocuous for a
country's internal ranking of sectors, and a cross-country comparison of the policy level must keep
the two bases apart: the mean Policy index is 0.64 for the economies scored on both components and
0.39 for those scored on the IEA database alone. The Strategic Index uses the Policy index only
through within-country ranks, so the difference in levels does not move one country's index
against another's; the difference in composition does change which sectors a country ranks first.
For the countries the Tracker rates, the within-country policy rank correlates with the technology
weight $g_t$ at a median Spearman of 0.82 among the 40 rated ``insufficient'' and $-0.92$ among the
nine rated ``critically insufficient''; for the countries on the IEA database alone the median is
0.12. Part of what the policy term contributes to the Strategic Index of a rated country is
therefore the Tracker's grade applied to the emission index that the climate term already carries.

\begin{figure}[tbp]
  \centering
  \includegraphics[width=\textwidth]{strategic_policy}
  \caption{The basis of the Policy index. Panel (a): number of economies by the components of
  their Policy index, for all economies with an index (215), the economies with at least one
  Strategic Index cell (174) and the coalition members (57). ``IEA PAMS + CAT'' combines the IEA
  Policies and Measures database (weight 3) with the Climate Action Tracker rating (weight 2),
  national or, for EU members, the rating of the Union; ``IEA PAMS only'' marks economies the
  Tracker does not rate, scored on the IEA database alone. Panel (b): Spearman rank correlations
  among the two components and the composite over the 33 headline technology--stage cells of the
  65 economies scored on both: the median, interquartile range and range of the within-country
  correlations (59 economies whose Tracker component varies across cells), the correlation pooled
  over cells and the correlation of country means. Source: \texttt{policy\_component\_tbl} and
  \texttt{policy\_index} in \texttt{index\_outputs.rds}.}
  \label{fig:strategic_policy}
\end{figure}

\paragraph{Sample of the Strategic Index.} The Strategic Index is defined for 3,154 cells of 174
countries in 30 sectors. A cell enters only if the country has a Policy index and both of its
pillar cells are rankable, that is, the economy is sovereign and at most half of the index weight
of each pillar cell is imputed. Of the 6,476 cells of policy-eligible countries, 1,000 are excluded
because the Economic Opportunity cell alone is not rankable, 721 because the Energy Security cell
alone is not rankable, 979 because neither is, and 622 because the Energy Security cell is missing;
a further 442 cells belong to countries without a Policy index. Geothermal Midstream is excluded by
rule, because its Energy Security pillar is constant, and no further sector is degenerate; of the
32 reported sectors, Electric Grid Downstream and Geothermal Upstream have no scored cell. The
sample is smaller than the 192 countries of the last intermediate run, because the new Upstream
trade baskets add an imputed Trade category to the Upstream cells of small economies and push their
imputed weight share above the threshold. Of the 174 countries, 65 are scored on both policy
components and 109 on the IEA database alone. Every coalition member has Strategic Index cells,
1,500 in all, with a median of 29 sectors per member (from 9 to 30).

\paragraph{Distribution and the dominant pattern.} The index has mean 0.535, standard deviation
0.082 and median 0.542, from 0.273 (Australia, Coal Upstream) to 0.761 (Togo, Nuclear Upstream).
It is organised by sector far more than by country. Sector fixed effects explain 72.6 percent of
its variance, technology fixed effects alone 67.5 percent, and country fixed effects 4.4 percent;
the median standard deviation of a sector's index across countries is 0.039. The ordering of
sectors follows the two fixed terms. Solar Upstream (mean 0.634), Batteries Upstream (0.631),
Nuclear Upstream (0.620), Batteries Midstream (0.618) and Batteries Downstream (0.606) rank highest
on average, and the three coal stages and Oil Upstream and Downstream lowest (0.410--0.441), the
order of the technology weights $g_t$, which run from 0 for oil and 0.02 for coal to 0.999 for
nuclear and wind. Upstream cells have the highest mean index (0.560, against 0.521 Midstream and
0.527 Downstream), because their mean security need $1-\ES$ is highest (0.677, against 0.513 and
0.552): most countries import the resources that are extracted. The highest cells in the sample
are Nuclear Upstream cells of small low-income economies (Togo, Burkina Faso, South Sudan,
Afghanistan, Sudan, Timor-Leste and the Dominican Republic, 0.748--0.761). They combine the largest
technology weight with an Energy Security cell at the floor of its sector: 86 of the 174 countries
scored in Nuclear Upstream share one value, 0.172, that of an economy without uranium resources or
production. Nuclear Upstream is the highest sector of 62 countries, 43 of them at that floor and 55
without any recorded uranium resources or production. This is a construction effect of the need
reading at a floor of the Energy Security pillar, not evidence of an opportunity, and it is the
reason for the endowment screen of the allied pick. Four other sectors have more than a quarter of
their countries at one Energy Security value: Green Hydrogen Midstream (34 of 40 at 0.256), Green
Hydrogen Downstream (45 of 158), Nuclear Downstream (19 of 67) and Batteries Downstream (38 of
138). No allied pick lies in a floor cell.

\begin{figure}[tbp]
  \centering
  \includegraphics[width=\textwidth]{strategic_heatmap}
  \caption{The Strategic Index across the coalition. Rows are the 57 coalition members (31 core
  members, then 26 partners), columns the 33 headline technology--stage cells grouped by stage;
  colour is $\SI_{its}$ (equation~\eqref{eq:si}). Grey cells are not scored (no rankable
  Economic Opportunity or Energy Security cell, a missing Energy Security cell, or Geothermal
  Midstream, excluded by rule). A dot marks an Upstream cell in which the member has no recorded
  reserves or production of the stage's resources: the endowment screen bars it from the allied
  pick, but its index value is unchanged. The outlined cell in each row is the member's allied
  pick (equation~\eqref{eq:si_allied}) with the screen. Source: \texttt{strategic\_index} in
  \texttt{index\_outputs.rds}; \texttt{allied\_strategic.csv}.}
  \label{fig:strategic_heatmap}
\end{figure}

\paragraph{Decomposition.} Averaged over the 167 countries with at least five cells, the five
weighted terms account for 18.8 percent (policy), 27.6 percent (Economic Opportunity), 32.4
percent (security need), 9.3 percent (stage weight) and 11.9 percent (technology weight) of the
index level. Their contributions to the within-country variance, measured as covariance shares
that sum to one, are very different: on average across countries, 40.7 percent for the technology
weight, 34.1 percent for the need term and 32.0 percent for the policy term, against $-1.4$ percent
for Economic Opportunity and $-5.4$ percent for the stage weight. The within-country Spearman
correlation of the index with each term has a median of 0.72 for the technology weight, 0.66 for
policy and 0.57 for need, but $-0.08$ for Economic Opportunity. Economic Opportunity carries the
second largest weight and almost none of the within-country variation because, within a country,
it moves against the need term: the median within-country Spearman correlation between $\EO$ and
$1-\ES$ is $-0.47$, negative in 164 of 167 countries. A country's strongest opportunities tend to
be the sectors in which it is already secure, so the two terms largely offset each other, and the
ranking of a country's sectors is set by the technology weight, the need for security and the
policy rank. The coalition gives the same picture (technology weight 40.2 percent, need 34.8
percent and policy 25.9 percent of the within-country variance; median
$\mathrm{corr}(\EO,1-\ES)=-0.49$, negative for every member). Figure~\ref{fig:strategic_decomp}
shows the five terms for the five highest sectors of six members. The Downstream sectors that head
the lists of Japan, Korea, Australia, Germany and India rank first on policy in those countries,
and most of the Upstream sectors that follow for Japan, Korea and Germany are sectors in which they
hold no resources, which the endowment screen bars from the pick.

\begin{figure}[tbp]
  \centering
  \includegraphics[width=\textwidth]{strategic_decomp}
  \caption{Decomposition of the Strategic Index for the five highest-ranked sectors of six
  coalition members. Each bar stacks the weighted terms of equation~\eqref{eq:si}:
  $0.2\,\POL^{*}_{its}$ (the within-country S-curve of the Policy index), $0.3\,\EO_{its}$,
  $0.3\,(1-\ES_{its})$, $0.1\,\lambda_s$ (Upstream 0.5, Midstream 0.75, Downstream 0.25) and
  $0.1\,g_t$; the bar length is $\SI_{its}$. (U), (M) and (D) mark the stage; an asterisk marks the
  member's allied pick and a circle an Upstream sector barred from the pick by the endowment
  screen. Source: \texttt{strategic\_index} in \texttt{index\_outputs.rds}.}
  \label{fig:strategic_decomp}
\end{figure}

\paragraph{Allied picks.} Every one of the 57 coalition members has an allied pick, the highest
Upstream or Midstream sector outside Geothermal and Green Hydrogen among the sectors it may pick,
and none has a tie (Table~\ref{tab:strategic_allied}). Batteries Midstream is the pick of 27
members (18 core members and nine partners, among them the United States, Germany, Korea, India,
Australia, France, the United Kingdom, Brazil and Mexico), Nuclear Midstream of nine (Costa Rica,
Czechia, the Dominican Republic, Guatemala, Hungary, Indonesia, Namibia, Qatar and Slovakia),
Nuclear Upstream of eight (Chile, the Democratic Republic of the Congo, Finland, Gabon, Greece,
Japan, Saudi Arabia and Sweden, all with identified uranium resources), Batteries Upstream of four
(Argentina, Norway, Peru and Ukraine), Electric Grid Midstream of four (Canada, Colombia, Morocco
and Oman), Solar Midstream of three (Israel, New Zealand and Singapore) and Wind Upstream of two
(Ghana and Poland); 43 picks are Midstream and 14 Upstream. The 11 members scored on the IEA
database alone all have picks, four of them in Nuclear Midstream. Within a member's eligible
sectors (a median of 11, from 2 to 18), the pick ranks first on the policy term for 22 members
and on the need term for 23, but first on Economic Opportunity for none and among the top three on
it for only seven: the picks are sectors a member needs and supports with policy rather than
sectors in which it is already competitive. Averaged over members, the pick exceeds its runner-up
mainly through the need term (0.011) and the technology weight (0.008).

The picks are not stable to small changes. The median margin between the pick and its runner-up is
0.015; it is below 0.005 for 10 members (17.5 percent: the United Kingdom, Hungary, Indonesia,
Luxembourg, Morocco, Oman, Portugal, Qatar, Ukraine and South Africa) and below 0.01 for 21 (36.8
percent), with a minimum of 0.0005 (Luxembourg). These margins are small against the index's
within-country standard deviation, 0.080 on average, so for a third of the members the pick is a
choice between two nearly equal sectors, and Table~\ref{tab:strategic_allied} reports the runner-up
beside it. The distribution of picks also moved during the final data revisions. In the run of 8
October, with the screen, Electric Grid Midstream was the pick of 17 members and Batteries
Midstream of 11; completing the energy-trade file filled the missing codes of the grid basket, and
the two sectors now account for four and 27 picks.

The Nuclear Midstream picks of the small partners call for a specific reading. Costa Rica, the
Dominican Republic and Guatemala each have four sectors they may pick, the Midstream stages of
nuclear, oil, gas and coal; Namibia has the same four and Nuclear Upstream. Their other Midstream
cells are not rankable, and their Upstream cells are either not rankable or barred by the screen.
Among these sectors nuclear has by far the largest technology weight ($g_t=0.999$, against 0.42 for
gas, 0.02 for coal and 0 for oil), which separates the pick from the runner-up by 0.058 (Costa Rica
and the Dominican Republic) and 0.100 (Guatemala), more than any other term. Nuclear Midstream
Energy Security is compressed around its median (interquartile range 0.095 across the 154
countries scored, against 0.12--0.15 in the fossil Midstream sectors) and somewhat lower (median
0.48, against 0.51--0.57), so the need term adds a smaller premium (0.020--0.042). These picks are
therefore not an Energy Security floor effect: they are what the technology weight selects when a
member's choice set holds only the fuel-processing stages. Without the screen, Costa Rica, the
Dominican Republic and Guatemala would pick Nuclear Upstream instead, at an index higher by 0.085
to 0.216, on Energy Security cells at the uranium floor.

% Generated by scripts/paper/fig_strategic.R -- do not edit by hand.
\begingroup
\singlespacing\scriptsize
\setlength{\tabcolsep}{2.9pt}
\setlength{\LTcapwidth}{\textwidth}
\begin{longtable}{@{}l l S[table-format=1.3] l S[table-format=1.4] c l c c@{}}
\caption{Allied strategic selection: pick, runner-up, margin and robustness of the pick}
\label{tab:strategic_allied}\\
\toprule
 & & & & & & & \multicolumn{2}{c@{}}{\makecell{Same pick\\under}}\\
\cmidrule(l){8-9}
Member & Pick & {$\SI$} & Runner-up & {Margin} & \makecell{AIB\\coalitions} & \makecell[l]{Pick without\\the screen} & (1) & (2)\\
\midrule
\endfirsthead
\multicolumn{9}{@{}l}{\small\textbf{Table \thetable:} continued}\\
\toprule
 & & & & & & & \multicolumn{2}{c@{}}{\makecell{Same pick\\under}}\\
\cmidrule(l){8-9}
Member & Pick & {$\SI$} & Runner-up & {Margin} & \makecell{AIB\\coalitions} & \makecell[l]{Pick without\\the screen} & (1) & (2)\\
\midrule
\endhead
\midrule
\multicolumn{9}{r@{}}{\textit{Continued on the next page}}\\
\endfoot
\bottomrule
\endlastfoot
\multicolumn{9}{@{}l}{\textit{Core members (31 of 31 with a pick)}}\\
Australia & Batteries (M) & 0.610 & Solar (M) & 0.0063 & 1 & same & -- & $\checkmark$ \\
Austria & Batteries (M) & 0.647 & Solar (M) & 0.0187 & 2 & Nuclear (U) & -- & $\checkmark$ \\
Belgium & Batteries (M) & 0.634 & Solar (M) & 0.0142 & -- & Wind (U) & -- & $\checkmark$ \\
Bulgaria & Batteries (M) & 0.651 & Nuclear (U) & 0.0101 & -- & Batteries (U) & -- & -- \\
Canada & Electric grid (M) & 0.615 & Batteries (M) & 0.0052 & -- & same & -- & -- \\
Croatia & Batteries (M) & 0.638 & Electric grid (M) & 0.0235 & -- & Nuclear (U) & -- & $\checkmark$ \\
Czech Republic & Nuclear (M) & 0.632 & Solar (M) & 0.0071 & -- & Wind (U) & -- & $\checkmark$ \\
Denmark & Batteries (M) & 0.636 & Solar (M) & 0.0109 & -- & Nuclear (U) & -- & $\checkmark$ \\
Finland & Nuclear (U) & 0.658 & Batteries (M) & 0.0245 & -- & same & -- & $\checkmark$ \\
France & Batteries (M) & 0.642 & Electric grid (M) & 0.0421 & 1 & Nuclear (U) & -- & $\checkmark$ \\
Germany & Batteries (M) & 0.623 & Nuclear (U) & 0.0073 & 2 & Solar (U) & -- & -- \\
Greece & Nuclear (U) & 0.643 & Batteries (M) & 0.0252 & -- & Batteries (U) & -- & $\checkmark$ \\
Hungary & Nuclear (M) & 0.606 & Wind (M) & 0.0015 & -- & Batteries (U) & -- & -- \\
Ireland & Batteries (M) & 0.636 & Solar (M) & 0.0244 & -- & Nuclear (U) & -- & $\checkmark$ \\
Israel$^{\dagger}$ & Solar (M) & 0.636 & Electric grid (M) & 0.0487 & -- & Batteries (U) & -- & $\checkmark$ \\
Italy & Batteries (M) & 0.639 & Nuclear (U) & 0.0181 & -- & Batteries (U) & -- & -- \\
Japan & Nuclear (U) & 0.620 & Batteries (M) & 0.0231 & -- & Batteries (U) & -- & $\checkmark$ \\
Luxembourg & Batteries (M) & 0.628 & Solar (M) & 0.0005 & -- & Nuclear (U) & -- & $\checkmark$ \\
Netherlands & Batteries (M) & 0.635 & Electric grid (M) & 0.0242 & 1 & Nuclear (U) & -- & $\checkmark$ \\
New Zealand & Solar (M) & 0.587 & Electric grid (M) & 0.0269 & -- & Batteries (U) & -- & $\checkmark$ \\
Norway & Batteries (U) & 0.618 & Batteries (M) & 0.0148 & -- & Solar (U) & -- & $\checkmark$ \\
Poland & Wind (U) & 0.633 & Batteries (U) & 0.0089 & -- & Nuclear (U) & -- & $\checkmark$ \\
Portugal & Batteries (M) & 0.644 & Nuclear (U) & 0.0007 & -- & Wind (U) & -- & -- \\
Romania & Batteries (M) & 0.643 & Nuclear (U) & 0.0195 & -- & Batteries (U) & -- & -- \\
Slovak Republic & Nuclear (M) & 0.644 & Batteries (M) & 0.0054 & -- & Wind (U) & -- & -- \\
South Korea & Batteries (M) & 0.635 & Batteries (U) & 0.0348 & 2 & same & -- & $\checkmark$ \\
Spain & Batteries (M) & 0.636 & Nuclear (U) & 0.0200 & -- & same & -- & -- \\
Sweden & Nuclear (U) & 0.649 & Wind (U) & 0.0145 & -- & Batteries (U) & -- & $\checkmark$ \\
Ukraine$^{\dagger}$ & Batteries (U) & 0.624 & Solar (M) & 0.0018 & -- & Solar (U) & -- & $\checkmark$ \\
United Kingdom & Batteries (M) & 0.617 & Electric grid (M) & 0.0042 & 1 & Nuclear (U) & -- & $\checkmark$ \\
United States & Batteries (M) & 0.632 & Electric grid (M) & 0.0096 & 3 & same & -- & $\checkmark$ \\
\midrule
\multicolumn{9}{@{}l}{\textit{Partners (26 of 26 with a pick)}}\\
Argentina & Batteries (U) & 0.551 & Solar (M) & 0.0114 & -- & Solar (U) & -- & -- \\
Brazil & Batteries (M) & 0.606 & Solar (U) & 0.0088 & -- & Electric grid (U) & -- & -- \\
Chile & Nuclear (U) & 0.645 & Batteries (M) & 0.0251 & -- & same & -- & $\checkmark$ \\
Colombia & Electric grid (M) & 0.597 & Solar (M) & 0.0195 & -- & Batteries (U) & $\checkmark$ & $\checkmark$ \\
Costa Rica & Nuclear (M) & 0.592 & Gas (M) & 0.1285 & -- & Nuclear (U) & $\checkmark$ & $\checkmark$ \\
Democratic Republic of Congo$^{\dagger}$ & Nuclear (U) & 0.664 & Batteries (U) & 0.0053 & -- & same & -- & $\checkmark$ \\
Dominican Republic$^{\dagger}$ & Nuclear (M) & 0.532 & Gas (M) & 0.0733 & -- & Nuclear (U) & $\checkmark$ & -- \\
Gabon$^{\dagger}$ & Nuclear (U) & 0.616 & Oil (U) & 0.0921 & -- & same & -- & -- \\
Ghana$^{\dagger}$ & Wind (U) & 0.589 & Electric vehicles (U) & 0.0469 & -- & same & -- & $\checkmark$ \\
Guatemala$^{\dagger}$ & Nuclear (M) & 0.546 & Oil (M) & 0.0768 & -- & Nuclear (U) & $\checkmark$ & -- \\
India & Batteries (M) & 0.628 & Wind (M) & 0.0437 & 1 & same & -- & $\checkmark$ \\
Indonesia & Nuclear (M) & 0.568 & Solar (M) & 0.0007 & 1 & same & -- & -- \\
Malaysia$^{\dagger}$ & Batteries (M) & 0.596 & Solar (M) & 0.0109 & 1 & Solar (U) & -- & $\checkmark$ \\
Mexico & Batteries (M) & 0.564 & Electric grid (M) & 0.0064 & -- & same & -- & -- \\
Morocco & Electric grid (M) & 0.609 & Batteries (M) & 0.0014 & -- & Solar (U) & $\checkmark$ & -- \\
Namibia$^{\dagger}$ & Nuclear (M) & 0.542 & Nuclear (U) & 0.0223 & -- & same & -- & -- \\
Oman$^{\dagger}$ & Electric grid (M) & 0.556 & Nuclear (M) & 0.0011 & -- & Solar (U) & -- & -- \\
Peru & Batteries (U) & 0.639 & Nuclear (U) & 0.0189 & -- & same & -- & $\checkmark$ \\
Philippines & Batteries (M) & 0.648 & Solar (U) & 0.0190 & -- & Nuclear (U) & -- & -- \\
Qatar$^{\dagger}$ & Nuclear (M) & 0.573 & Oil (M) & 0.0031 & -- & Batteries (U) & -- & -- \\
Saudi Arabia & Nuclear (U) & 0.567 & Electric grid (M) & 0.0156 & -- & Solar (U) & -- & -- \\
Singapore & Solar (M) & 0.601 & Electric grid (M) & 0.0084 & -- & Solar (U) & -- & -- \\
South Africa & Batteries (M) & 0.620 & Solar (M) & 0.0019 & 1 & Solar (U) & -- & $\checkmark$ \\
Thailand & Batteries (M) & 0.587 & Solar (M) & 0.0182 & -- & Solar (U) & -- & -- \\
Turkiye & Batteries (M) & 0.556 & Solar (M) & 0.0124 & -- & Solar (U) & -- & -- \\
Viet Nam & Batteries (M) & 0.606 & Nuclear (U) & 0.0273 & 2 & same & -- & -- \\
\midrule
\multicolumn{6}{@{}l}{Members with the same pick (of 57)} & 16 & 5 & 31\\
\end{longtable}
\vspace{-6pt}\noindent
\textit{Notes:} For each coalition member with a Strategic Index, the pick is the Upstream or Midstream sector outside Geothermal and Green Hydrogen with the highest index (equation~\eqref{eq:si_allied}), after the endowment screen bars Upstream sectors in which the member has no recorded reserves or production; the runner-up is the second-highest such sector and the margin the difference between their indices; (U) Upstream, (M) Midstream. No member has a tie. AIB coalitions counts the Sectoral Compacts of Section~\ref{sec:network} with a node in the pick stage whose exact minimum coalition includes the member (-- if none). ``Pick without the screen'' is the pick when the endowment screen is switched off (``same'' if unchanged); the index values are the same in both cases. Columns (1) and (2) mark whether the pick, with the screen, is unchanged when the index is recomputed on the same cells (1) with security strength $\ES$ in place of need $1-\ES$ and (2) without the stage and technology weights. All 57 members have a pick. $^{\dagger}$\,Policy index on IEA PAMS alone (no CAT rating; 11 members); the other members' Policy index combines IEA PAMS and CAT. Source: \texttt{index\_outputs.rds} and \texttt{allied\_strategic.csv} of the final run, which the script reproduces; \texttt{scripts/paper/fig\_strategic.R}.\par
\endgroup
\medskip


\paragraph{The endowment screen.} The screen makes an Upstream extraction stage eligible for a
member's pick only if the member has recorded reserves or production of the stage's resources; it
leaves every index value unchanged. Of the 998 scored Upstream cells in the ten screened
technologies, 669 have no recorded resource base, 306 of the coalition's 509. The screen changes
the pick of 41 of the 57 members (Figure~\ref{fig:strategic_screen}). Without it, 47 picks would be
Upstream (Nuclear Upstream 17, Batteries Upstream 12, Solar Upstream 12 and Wind Upstream 5), many
of them in members without the resource, such as Nuclear Upstream for France, the United Kingdom,
the Netherlands, Denmark, Ireland and Luxembourg and Solar Upstream for Germany, Saudi Arabia and
T\"urkiye. With it, 43 picks are Midstream, and the eight Nuclear Upstream picks that remain all
rest on identified uranium resources. The index forgone by the switch is small for most members
(median 0.026, interquartile range 0.016--0.040) and largest for the Dominican Republic (0.216) and
Colombia (0.128). Seven members move to the Midstream stage of the same technology (three each in
batteries and nuclear, one in solar). The screen is a rule of the pick, not of the index: it encodes
the judgement that an economy cannot become an allied supplier of a mined resource it does not
have, whatever its need for it.

\begin{figure}[tbp]
  \centering
  \includegraphics[width=\textwidth]{strategic_screen}
  \caption{Allied picks with and without the endowment screen. Panel (a): number of the 57
  coalition members whose pick is each sector, with the screen (filled markers, the headline) and
  without it (open markers); colour is the stage. Panel (b): the transition between the two
  selections; each tile counts the members whose pick without the screen is the row and with the
  screen the column, and outlined tiles on the diagonal are unchanged picks (16 members). The
  screen bars an Upstream stage from a member's pick unless the member has recorded reserves or
  production of the stage's resources; index values are the same in both selections, and the
  screened selection equals \texttt{allied\_strategic.csv}. Source: \texttt{strategic\_index} in
  \texttt{index\_outputs.rds}; selections recomputed with
  \texttt{select\_allied\_strategic\_sectors()}, \texttt{require\_resource\_base} TRUE and FALSE.}
  \label{fig:strategic_screen}
\end{figure}

\paragraph{Robustness to the security sign and to the other normative choices.} The Strategic
Index reads security as need: a sector scores higher where the country is less secure. Reading it
as strength, with $\ES$ in place of $1-\ES$ and every other term unchanged, gives an index that
agrees only weakly with the headline (Figure~\ref{fig:strategic_sign}). Within members the median
Spearman correlation of the two readings is 0.36 (interquartile range 0.17--0.48); it is 0.38 for
core members and 0.24 for partners, and it is negative for ten partners (Argentina, Indonesia,
Mexico, Namibia, Oman, Qatar, Saudi Arabia, Thailand, T\"urkiye and Viet Nam), whose most secure
sectors are their fuel and processing industries. Only 5 of the 57 picks survive the change of
sign, and for 12 members the headline pick is among the top three of the strength reading, whose
picks move to Nuclear Midstream (28), Gas Midstream (10) and Electric Grid Midstream (7). Without
the stage and technology weights the two readings are unrelated (median within-member correlation
0.01, one pick in common). The sign of the security term is therefore the most consequential
normative choice in the index, and it is a choice: the need reading asks where an economy should
build capability to reduce its exposure, the strength reading where it can build on capability it
has.

The other choices matter less for the ordering of a member's sectors but more for its single top
pick. Dropping the stage and technology weights keeps the need reading close to the headline
(median within-member Spearman 0.90) but keeps only 31 of the 57 picks, although the headline pick
stays among the variant's top three for 49 members. Ranking Economic Opportunity and need within
the country, like the policy term, gives a median of 0.89, 20 identical picks and 49 in the top
three. Dropping the policy term gives 0.88, 30 identical picks and 55 in the top three, and moves
picks toward Nuclear Upstream (20) and Electric Grid Midstream (13). For the 46 members scored on
both policy components, building the policy term from the IEA database alone gives 0.89 and 34
identical picks. Adding back the industrial-policy measures of the New Industrial Policy
Observatory \citep{gta2026nipo}, from the October 2026 run (77 jurisdictions, measures in force on
5 October 2026), with the weights 4, 3 and 2 that the Observatory, the IEA component and the
Tracker carried before October 2026, gives 0.96 for the 50 members with Observatory data, 27
identical picks and 44 in the top three. Requiring both policy components, the rule before October
2026, leaves 65 eligible countries and 1,666 Strategic Index cells whose values and picks are
identical to the headline; it only removes the 11 members scored on the IEA database alone. The
ranking of sectors within a member is thus robust to the treatment of policy and of the fixed
weights over the whole list, but the top pick is not: it changes for between a quarter and two
thirds of the members under each of these variants, as the small margins imply.

\begin{figure}[tbp]
  \centering
  \includegraphics[width=\textwidth]{strategic_sign}
  \caption{Sensitivity of the Strategic Index to its normative choices, coalition members. Panel
  (a): within-member Spearman rank correlation between the headline index, which reads security as
  need ($1-\ES$), and the index with security as strength ($\ES$), all other terms unchanged; one
  point per member, by tier, with the median marked. Panel (b): within-member Spearman correlation
  of each variant with the headline (median, interquartile range and range): security as strength;
  strength without the stage and technology weights; need without those weights; Economic
  Opportunity and need ranked within the country; no policy term; the policy term from the IEA
  Policies and Measures database alone (46 members scored on both components); the October 2026
  component of the New Industrial Policy Observatory added back to the policy composite (50
  members with Observatory data). Panel (c): share of members whose allied pick is unchanged
  (filled) and whose headline pick is among the variant's top three sectors (open). All variants
  keep the endowment screen and the cells of the headline index. Source:
  \texttt{index\_outputs.rds}; the Observatory component of the October 2026 run.}
  \label{fig:strategic_sign}
\end{figure}
